%% RNAAS research note: injection-recovery sensitivity of a TESS 20-s transit search of white dwarfs within 100 pc
%% Compile with AASTeX 6.3.1 (e.g. on Overleaf: upload this file, fig1.pdf and note.bib).
%% Numbers: data/inject/analysis.json and data/inject/injections.csv (python -m wdplanets.inject analyse / plot).
\documentclass[RNAAS]{aastex631}

\begin{document}

\title{How Small a Transiting Body Can TESS 20-s Photometry Detect around Nearby White Dwarfs? An Injection--Recovery Test}

\author{Christian Jair}
\affiliation{Independent researcher}

\begin{abstract}
We measure the sensitivity of a box-least-squares transit search of TESS 20-s light curves of white dwarfs within
100\,pc by injecting 304 physical transit models (0.1--10\,$R_\oplus$, periods 0.2--10\,d, all transiting
orientations) into the raw light curves of 16 already-searched white dwarfs ($G=13.6$--16.7) and running the
unmodified search and thresholds. 57\% of the bodies are recovered at the correct period and none at a wrong one;
every non-grazing body of $\geq1\,R_\oplus$ is recovered. The radius recovered half of the time at $P=1$\,d is 0.26,
0.41, 0.67 and 1.08\,$R_\oplus$ at $G\approx13.5$, 14.5, 15.5 and 16.5, and scales with magnitude, period and number of
sectors as expected for a background-limited box search. Applied to the 1,489 white dwarfs searched so far, an
Earth-sized body at 0.2--10\,d would have been detected with $\geq50\%$ probability around 82\% of them, a Moon-sized
body around only 4.5\%.
\end{abstract}

\keywords{White dwarf stars (1799) --- Transit photometry (1709) --- Exoplanet detection methods (489) --- Exoplanets (498)}

\section{Introduction}
A white dwarf is about the size of the Earth, so a transiting planet blocks a large fraction of its light, but for
only a few minutes. WD\,1856+534\,b, a giant planet on a 1.4-d orbit, is the only planet known to transit one
\citep{vanderburg2020}. A survey that finds nothing is informative only if its sensitivity is known. For K2,
\citet{vansluijs2018} found that bodies larger than 0.5\,$R_\oplus$ (1-min cadence) or 1\,$R_\oplus$ (30-min cadence)
were recovered more often than not at periods below 16\,h. \citet{robert2024} searched TESS 120-s and 20-s data of
313 polluted white dwarfs and presented their sensitivity as period--radius maps. Here we report the sensitivity of
a search of TESS 20-s data of the white dwarfs within 100\,pc as a function of radius, brightness and period, and
what it implies for the stars searched so far.

\section{Search and injections}
The search targets white dwarfs within 100\,pc in the \citet{gentilefusillo2021} catalogue with $G<17$ and SPOC light
curves, using the six newest sectors at 20-s cadence (120\,s where 20-s data do not exist). Each PDCSAP light curve is
flattened over 12\,h, clipped upward only, and binned to 40\,s; a box-least-squares search \citep{kovacs2002} with
boxes of 1--30\,min is run per sector on a shared grid of 176,000 periods between 0.15 and 13\,d, the per-sector
log-likelihoods are summed, and the best period is refined on all sectors. A signal is flagged when its signal
detection efficiency is $\geq10$, its depth S/N $\geq8$ and it has $\geq3$ events.

We chose 12 test stars among those already searched with no signal above threshold (three per $G$ bin: $<14$,
14--15, 15--16, 16--17), each with 20-s data in 2--6 sectors, and injected 25 bodies into each; four further stars
carry one injection each from an earlier trial. The white-dwarf radius follows from the catalogue's mass and
$\log g$; the orbit is circular, with $a$ from Kepler's third law. The model is the exact overlap of two uniform disks
(so grazing and larger-than-the-star bodies are handled), averaged over each exposure. Radii and periods are drawn
log-uniformly in 0.1--10\,$R_\oplus$ and 0.2--10\,d, the impact parameter uniformly in $0$--$(1+R_p/R_\star)$, and the
phase uniformly. The model multiplies the raw normalized flux \emph{before} detrending, and the unmodified search runs
on the result. A body is recovered if it is flagged and its period matches the injected one, or a $\times\frac12$,
2, $\frac13$ or 3 alias, within 0.2\%.

We fit the 304 outcomes with $p = A/\{1+\exp[-(\log_{10}R_p-\mu)/s]\}$, where
$\mu=\mu_0+c_G(G-15)+c_P\log_{10}(P/{\rm d})+c_N\log_{10}N_{\rm sec}$ and the ceiling $A$ accounts for grazing
orbits; uncertainties come from 500 bootstrap refits resampling stars.

\section{Results}
Of the 304 bodies, 174 (57.2\%) were recovered, and all 174 flagged signals were at the correct period
(Figure~\ref{fig:1}). Recovery is 11\%, 44\%, 78\% and 94\% for 0.1--0.3, 0.3--1, 1--3 and 3--10\,$R_\oplus$. All 65
non-grazing bodies with $R_p\geq1\,R_\oplus$ were recovered; the 21 misses in that range are grazing orbits that hid at
most 19\% of the disk, giving $A=0.91$ (0.86--0.94). The fitted coefficients are $c_G=0.21$ (0.16--0.26)\,dex per
magnitude, $c_P=0.21$ (0.11--0.26) and $c_N=-0.28$ ($-0.46$ to $-0.06$), with a transition width $s=0.09$\,dex. A box
search whose S/N scales as $(R_p/R_\star)^2\sqrt{N_{\rm in}}/\sigma$, with the number of in-transit points
$N_{\rm in}\propto N_{\rm sec}P^{-2/3}$, predicts $c_P=1/6$ and $c_N=-1/4$, and $c_G=0.2$ when the noise is
background-limited. The radius recovered half of the time at $P=1$\,d and four sectors is 0.26 (0.21--0.31), 0.41
(0.34--0.48), 0.67 (0.52--0.80) and 1.08 (0.76--1.42)\,$R_\oplus$ at $G=13.5$, 14.5, 15.5 and 16.5; at $G=15$ it grows
from 0.38\,$R_\oplus$ at 0.2\,d to 0.85\,$R_\oplus$ at 10\,d. Individual stars recovered between 36\% and 80\% of the
same distribution of bodies, the lowest being faint stars or stars with two sectors.

We applied the model to each of the 1,489 white dwarfs searched so far ($G=11.0$--17.0, median four sectors),
averaging over a log-uniform period distribution in 0.2--10\,d. An Earth-sized body would have been detected with
$\geq50\%$ probability around 82\% (68--97\%) of them, and with a mean probability of 69\% (60--82\%); for a Moon-sized
body (0.27\,$R_\oplus$) the figures are 4.5\% (3--8\%) and 8\% (6--12\%). Multiplying by each star's transit
probability (median 0.6\% for an Earth at 1\,d) gives an effective sample of 7.5 stars for Earth-sized bodies; if no
such body is confirmed among the remaining candidates, the 95\% upper limit on their occurrence at 0.2--10\,d would
be $\approx30\%$ \citep[eq.~7 of][]{vansluijs2018}.

\section{Discussion}
Twenty-second TESS photometry reaches sub-Earth radii out to several days around the brightest nearby white dwarfs,
where the K2 search reached 0.5\,$R_\oplus$ only below 16\,h \citep{vansluijs2018}, but most of the 100\,pc sample is
fainter than $G=15$, where the limit is two thirds of an Earth radius or more, and Moon-sized bodies reach 50\% recovery
(at 1\,d) only around stars brighter than $G\approx13.6$. Because transit probabilities are below 1\%, even 1,500 stars
constrain the occurrence of Earth-sized bodies only to tens of percent, comparable with the K2 limit of $<28\%$ for
habitable-zone Earths \citep{vansluijs2018} and the TESS limits for polluted white dwarfs \citep{robert2024}; the
period ranges and samples differ. Our injections precede detrending, use each star's own mass and radius, and require
a 0.2\% period match, choices that, if anything, make the sensitivity conservative. The main limitations are the
small number of test stars, which dominates the uncertainty; the extrapolation to stars outside $G=13.6$--16.7, to
single-sector stars and to stars with only 120-s data; and the omission of the vetting steps, of limb darkening and
of non-periodic debris transits. The test stars had no prior signal, so the absence of wrong-period detections does
not measure the false-alarm rate on noisy stars.

\begin{figure}[ht!]
\plotone{fig1.pdf}
\caption{(a) Recovered/injected bodies by injected radius and Gaia $G$. (b) Recovered fraction versus radius in four
$G$ bins (points, 68\% Wilson intervals, averaged over period) with the fitted model at $P=1$\,d and four sectors
(lines). (c) All injected bodies in the period--radius plane: recovered (filled), missed (open), missed on a grazing
orbit (crosses); lines show the model's 50\% recovery radius at $G=13.5$, 14.5, 15.5 and 16.5.\label{fig:1}}
\end{figure}

\begin{acknowledgments}
This paper includes data collected by the TESS mission, obtained from MAST; funding for the TESS mission is provided by
the NASA Science Mission Directorate. It has made use of data from the European Space Agency (ESA) mission Gaia,
processed by the Gaia Data Processing and Analysis Consortium, and of NASA's Astrophysics Data System.
We made use of AI-assisted tools (Anthropic's Claude) to speed up data retrieval, code development and the
preparation of this note; all analyses were verified by the author, who is responsible for the content.
\end{acknowledgments}

\facilities{TESS, Gaia}
\software{astropy, astroquery, lightkurve, matplotlib, numpy, pandas, scipy}

\bibliography{note}{}
\bibliographystyle{aasjournal}

\end{document}
